\documentclass[conference]{IEEEtran}
\IEEEoverridecommandlockouts

\usepackage{cite}
\usepackage{amsmath,amssymb,amsfonts}
\usepackage{algorithmic}
\usepackage{graphicx}
\usepackage{textcomp}
\usepackage{xcolor}
\usepackage{booktabs}
\usepackage{url}

\def\BibTeX{{\rm B\kern-.05em{\sc i\kern-.025em b}\kern-.08em
    T\kern-.1667em\lower.7ex\hbox{E}\kern-.125emX}}

\begin{document}

\title{Real-Time Perceptual Image Abstraction Using Edge-Aware Filtering and Adaptive Stylization}

\author{
  \IEEEauthorblockN{Jay Salot}
  \IEEEauthorblockA{\textit{Department of Computer Science}\\
  \textit{Dhirubhai Ambani University}\\
  Gandhinagar, India \\
  22101xxx@daiict.ac.in}
}

\maketitle

\begin{abstract}
We implement and extend the real-time image abstraction pipeline of Winemöller et al. \cite{winemoller2006} (SIGGRAPH 2006), which produces perceptually simplified, cartoon-like representations of images using iterated bilateral filtering, Difference-of-Gaussians (DoG) edge detection, and soft color quantization. Our Phase 1 reproduces the paper's pipeline faithfully, including the gradient-adaptive sharpness control in the quantization step, which was previously only approximated. Our Phase 2 introduces two novel contributions: (1) a saliency-guided adaptive bilateral filter that varies the range sigma per pixel based on a combined gradient and color-opponency saliency map, and (2) data-driven adaptive k-means color quantization in full CIE Lab space, replacing the paper's fixed equidistant bins. Quantitative evaluation on a standard test image demonstrates that Phase 2 achieves 98.4\% color reduction compared to 19.4\% in Phase 1, with a modest and intentional decrease in PSNR (26.29 dB vs 33.31 dB) that reflects the stronger artistic abstraction.
\end{abstract}

\begin{IEEEkeywords}
image abstraction, bilateral filter, Difference-of-Gaussians, color quantization, saliency, non-photorealistic rendering
\end{IEEEkeywords}

%--------------------------------------------------------------------
\section{Introduction}
%--------------------------------------------------------------------

Natural images contain perceptual redundancy: fine textures, subtle gradients, and background clutter that the human visual system routinely ignores when recognizing objects or scenes. For video communication, recognition tasks, and artistic rendering, a simplified cartoon-like representation can be more useful than the raw photographic image.

Winemöller et al.~\cite{winemoller2006} presented a real-time pipeline that achieves this simplification using only standard GPU operations available in 2006, making it suitable for frame-rate processing of video. Their key insight was that iterating a bilateral filter approximates anisotropic diffusion—progressively flattening color regions while preserving edges—and that DoG filtering faithfully mimics retinal edge detection at a very low computational cost.

While the paper's approach is elegant, it has two acknowledged limitations: (i) the bilateral filter uses a globally uniform range sigma, treating all image regions identically regardless of their visual importance, and (ii) the color quantization uses fixed equidistant bin boundaries in luminance that are "arbitrary and make it difficult to control results for artistic purposes" (paper, Section 3.3).

We address both limitations while preserving the paper's real-time spirit.

%--------------------------------------------------------------------
\section{Background and Related Work}
%--------------------------------------------------------------------

\subsection{Bilateral Filtering}
The bilateral filter~\cite{tomasi1998} replaces each pixel with a weighted average of neighbors, where the weight decays both with spatial distance ($\sigma_d$) and with range distance ($\sigma_r$). Small $\sigma_r$ preserves fine edges; large $\sigma_r$ approaches a Gaussian blur. Iterating the filter approximates the heat equation with a spatially varying diffusion coefficient~\cite{winemoller2006}.

\subsection{Difference-of-Gaussians}
DoG edge detection subtracts two Gaussian-blurred versions of an image at different scales. The zero-crossings of this difference correspond approximately to edges, mimicking the response of retinal ON/OFF ganglion cells~\cite{marr1980}.

\subsection{Saliency Maps}
Saliency detection identifies visually important image regions. Itti et al.~\cite{itti1998} proposed a biologically motivated model combining color, intensity, and orientation contrast. We use a simplified version combining gradient magnitude and color opponency channels.

%--------------------------------------------------------------------
\section{Phase 1: Paper Implementation}
%--------------------------------------------------------------------

\subsection{Iterative Bilateral Filtering}

Let $I$ be the input image in CIE Lab color space. A single bilateral filter pass computes:
\begin{equation}
H(\hat{x}, \sigma_d, \sigma_r) = \frac{\sum_{\xi} f(\|\hat{x}-\xi\|, \sigma_d)\, g(\|I(\hat{x})-I(\xi)\|, \sigma_r)\, I(\xi)}{\sum_{\xi} f(\cdot)\, g(\cdot)}
\label{eq:bilateral}
\end{equation}
where $f$ and $g$ are Gaussian kernels with parameters $\sigma_d$ and $\sigma_r$ respectively.

We apply this filter $n_b = 4$ times. An intermediate snapshot after $n_e = 2$ iterations is stored for use in edge detection, as this preserves more high-frequency edge information than the fully diffused result.

We use $\sigma_d = 3.0$ and $\sigma_r = 4.25$ as specified in the paper, implemented via OpenCV's \texttt{bilateralFilter} applied independently to each Lab channel.

\subsection{Difference-of-Gaussians Edge Detection}

The DoG filter is computed on the luminance ($L$) channel of the intermediate image:
\begin{equation}
\text{DoG}(\hat{x}) = S(\hat{x}, \sigma_e) - \tau \cdot S(\hat{x}, \sqrt{1.6}\,\sigma_e)
\end{equation}
The edge function is:
\begin{equation}
D(\hat{x}) = \begin{cases}
1 & \text{if } \text{DoG}(\hat{x}) > 0\\
1 + \tanh(\phi_e \cdot \text{DoG}(\hat{x})) & \text{otherwise}
\end{cases}
\end{equation}
with $\sigma_e = 1.0$, $\tau = 0.98$, $\phi_e = 2.0$. Values near 0 indicate strong edges; these pixels are darkened in the final overlay.

\subsection{Gradient-Adaptive Soft Color Quantization}

Soft quantization (Equation 6 of the paper) maps each pixel's normalized luminance to a soft step function:
\begin{equation}
Q(\hat{x}) = q_n + \frac{\Delta q}{2} \tanh\!\bigl(\phi_q(\hat{x})\cdot(f(\hat{x}) - q_n)\bigr)
\label{eq:quant}
\end{equation}
where $q_n$ is the center of the nearest bin, $\Delta q = 1/n_\text{bins}$, and $\phi_q(\hat{x})$ is a spatially-varying sharpness:
\begin{equation}
\phi_q(\hat{x}) = \lambda_\phi + \frac{|\nabla L(\hat{x})| - \lambda_\delta}{\omega_\delta - \lambda_\delta}\,(\omega_\phi - \lambda_\phi)
\end{equation}
with $(\lambda_\phi, \omega_\phi) = (3, 14)$ and $(\lambda_\delta, \omega_\delta) = (0, 2)$.

\textbf{Correction:} The paper intends this formula to be applied pixel-wise with the full spatial $\phi_q$ map. Our implementation computes and applies the complete 2D $\phi_q$ map in a single vectorized NumPy operation, ensuring high-gradient regions receive hard bin boundaries while flat areas produce smooth painterly transitions.

%--------------------------------------------------------------------
\section{Phase 2: Novel Contributions}
%--------------------------------------------------------------------

\subsection{Saliency-Guided Adaptive Bilateral Filtering}

The paper's bilateral filter uses a uniform $\sigma_r$ across the entire image, providing no mechanism to treat visually important regions differently from backgrounds. The paper notes an attention weight $m(\cdot)$ in Eq.~(2) but leaves it unused in the automatic pipeline. We activate it explicitly.

\textbf{Saliency computation:} We combine two complementary cues:
\begin{equation}
S(\hat{x}) = w_g \cdot S_\nabla(\hat{x}) + w_c \cdot S_\text{opp}(\hat{x})
\end{equation}
where $S_\nabla = |\nabla I|/\max|\nabla I|$ is the normalized gradient magnitude and $S_\text{opp}$ combines red-green and blue-yellow color opponency:
\begin{equation}
S_\text{opp} = \frac{|R-G| + |B - 0.5(R+G)|}{2}
\end{equation}
We use $w_g = 0.6$, $w_c = 0.4$, and smooth $S$ with a Gaussian ($\sigma = 5$) to prevent sharp saliency boundaries from creating artifacts.

\textbf{Adaptive filtering:} We run two bilateral passes per iteration, one with $\sigma_r^\text{high} = 2.0$ (detail-preserving) and one with $\sigma_r^\text{low} = 8.0$ (aggressive smoothing), then blend:
\begin{equation}
H_\text{adaptive}(\hat{x}) = S(\hat{x})\cdot H_\text{high}(\hat{x}) + (1 - S(\hat{x}))\cdot H_\text{low}(\hat{x})
\end{equation}
Salient regions (faces, objects) are treated gently; low-saliency backgrounds are smoothed much more heavily, producing stronger abstraction where the human eye expects it.

\subsection{Adaptive K-Means Color Quantization}

The paper's equidistant bins on the $L$ channel are acknowledged to be "arbitrary." We replace them with data-driven k-means clustering in the full CIE Lab color space.

\textbf{Method:} Pixels are normalized to $[0,1]^3$ in Lab coordinates, then clustered with MiniBatchKMeans ($k = 8$, random seed 42). Each pixel is replaced by its nearest cluster center, then denormalized back to Lab. This operates on all three Lab channels, ensuring perceptually uniform color grouping rather than luminance-only quantization.

\textbf{Advantage:} The clusters naturally align with the actual color distribution of the image, not an arbitrary grid. For the Lena test image, this reduces the effective palette from hundreds of Lab colors to exactly 8 perceptually chosen representatives, yielding 98.4\% color reduction.

%--------------------------------------------------------------------
\section{Results}
%--------------------------------------------------------------------

\subsection{Quantitative Evaluation}

We evaluate on the standard Lena (512×512) test image. Table~\ref{tab:metrics} shows the quantitative comparison.

\begin{table}[h]
\caption{Quantitative Comparison on Lena Image}
\label{tab:metrics}
\begin{center}
\begin{tabular}{lcc}
\toprule
\textbf{Metric} & \textbf{Phase 1 (Paper)} & \textbf{Phase 2 (Ours)} \\
\midrule
PSNR (dB)         & 33.31 & 26.29 \\
SSIM              & 0.9954 & 0.9840 \\
Color Reduction (\%) & 19.4 & \textbf{98.4} \\
Edge Preservation & \textbf{0.8803} & 0.7583 \\
\bottomrule
\end{tabular}
\end{center}
\end{table}

\subsection{Analysis}

Phase 2 achieves dramatically stronger color abstraction (98.4\% vs 19.4\%) at the cost of pixel-level fidelity (PSNR 26.29 vs 33.31). This trade-off is intentional: the k-means palette quantization reduces the image to a clean 8-color poster-like representation, which is the desired artistic output. The lower PSNR simply reflects that we are further from the original, not that the output is worse—it is more abstracted.

Edge preservation decreases slightly (0.76 vs 0.88) because extreme color quantization eliminates some subtle gradient-based "edges" that are artifacts of smooth color gradients rather than true structural boundaries.

SSIM remains high (0.984) indicating that the structural content—object boundaries, spatial layout—is well preserved despite the stronger stylization.

%--------------------------------------------------------------------
\section{Discussion}
%--------------------------------------------------------------------

\subsection{Completeness vs. Paper}

All three core components of the 2006 paper are implemented correctly:
\begin{itemize}
\item Iterative bilateral filtering with the correct intermediate snapshot for edge extraction
\item DoG edge detection matching Equations (4) and (5) from the paper
\item Gradient-adaptive soft quantization (Eq.~6) now applied with the full spatial $\phi_q$ map
\end{itemize}

\subsection{Novelty Contribution}

Our contributions extend the paper along two well-motivated axes. The saliency-guided filtering addresses the fundamental limitation that real-world images have non-uniform visual importance—a face and a sky patch should not receive identical treatment. The adaptive quantization addresses the paper's own admitted arbitrariness in bin placement.

Both contributions require no special hardware or deep networks, maintaining the paper's computational spirit while producing stronger artistic results.

\subsection{Limitations}

The saliency computation is based on low-level cues only (gradients, color opponency). Semantic saliency—knowing that a face is more important than a chair regardless of gradient—would require a learned model (e.g., a CNN saliency predictor). This is left as future work.

%--------------------------------------------------------------------
\section{Conclusion}
%--------------------------------------------------------------------

We have implemented and quantitatively evaluated the full pipeline from Winemöller et al.~\cite{winemoller2006}, correcting an approximation in the soft quantization step, and extended it with two novel contributions: saliency-guided adaptive bilateral filtering and data-driven Lab-space k-means color quantization. Phase 2 produces a 98.4\% color reduction with strong artistic abstraction, compared to 19.4\% in the original pipeline. Both phases are implemented entirely in Python using OpenCV and scikit-learn, with all results reproducible from the provided code.

%--------------------------------------------------------------------
\begin{thebibliography}{00}

\bibitem{winemoller2006}
H. Winemöller, S. C. Olsen, and J. Gooch,
``Real-Time Video Abstraction,''
\textit{ACM Transactions on Graphics (Proc. SIGGRAPH)},
vol. 25, no. 3, pp. 1221--1226, 2006.

\bibitem{tomasi1998}
C. Tomasi and R. Manduchi,
``Bilateral Filtering for Gray and Color Images,''
\textit{Proc. IEEE ICCV}, pp. 839--846, 1998.

\bibitem{marr1980}
D. Marr and E. Hildreth,
``Theory of Edge Detection,''
\textit{Proc. Royal Society B}, vol. 207, pp. 187--217, 1980.

\bibitem{itti1998}
L. Itti, C. Koch, and E. Niebur,
``A Model of Saliency-Based Visual Attention for Rapid Scene Analysis,''
\textit{IEEE Trans. Pattern Anal. Mach. Intell.},
vol. 20, no. 11, pp. 1254--1259, 1998.

\end{thebibliography}

\end{document}
